\documentclass[9pt,a4paper]{extarticle}
\usepackage{parskip}

% =====================
% Codificação, idioma e layout
% =====================
\usepackage[T1]{fontenc}
\usepackage[utf8]{inputenc}
\usepackage[brazil]{babel}
\usepackage{lmodern}
\usepackage{geometry}
\geometry{margin=2cm}
\usepackage{microtype}

% =====================
% Matemática e símbolos
% =====================
\usepackage{amsmath, amssymb, amsthm}
\usepackage{mathtools}
\usepackage{bm}

% =====================
% Links
% =====================
\usepackage{hyperref}
\hypersetup{colorlinks=true,linkcolor=black,citecolor=black,urlcolor=blue}

% =====================
% Código-fonte (listings)
% =====================
\usepackage{listings}
\usepackage{xcolor}
\definecolor{codebg}{rgb}{0.96,0.96,0.96}
\definecolor{codekw}{RGB}{33,96,171}
\definecolor{codecmt}{RGB}{120,120,120}
\definecolor{codestr}{RGB}{163,21,21}
\lstdefinestyle{pystyle}{
  language=Python,
  backgroundcolor=\color{codebg},
  basicstyle=\ttfamily\footnotesize,
  keywordstyle=\color{codekw}\bfseries,
  commentstyle=\color{codecmt}\itshape,
  stringstyle=\color{codestr},
  showstringspaces=false,
  frame=single,
  breaklines=true,
  tabsize=2,
  inputencoding=utf8,
  extendedchars=true,
  literate=
   {á}{{\'a}}1 {à}{{\`a}}1 {ã}{{\~a}}1 {â}{{\^a}}1
   {é}{{\'e}}1 {ê}{{\^e}}1
   {í}{{\'i}}1
   {ó}{{\'o}}1 {õ}{{\~o}}1 {ô}{{\^o}}1
   {ú}{{\'u}}1
   {ç}{{\c{c}}}1
   {Á}{{\'A}}1 {Ã}{{\~A}}1
   {É}{{\'E}}1 {Ê}{{\^E}}1
   {Í}{{\'I}}1
   {Ó}{{\'O}}1 {Õ}{{\~O}}1 {Ô}{{\^O}}1
   {Ú}{{\'U}}1
   {Ç}{{\c{C}}}1
}
\lstset{style=pystyle}

% =====================
% Título
% =====================
\title{\textbf{TP — Aprendizado Contínuo com Controle por Reforço}\\\large Projeto Independente de Pesquisa}
\author{Entrega: Código, Relatório curto, Slides}
\date{\today}

\begin{document}
\maketitle

% ======================================================
% VISÃO GERAL
% ======================================================
\section{Visão geral}
O \textbf{aprendizado contínuo} (ou \emph{continual learning}) busca construir modelos que aprendem sequencialmente, incorporando novas tarefas sem esquecer as anteriores. O principal desafio é o \textit{catastrophic forgetting}, isto é, a perda de desempenho em tarefas já aprendidas ao adquirir novos conhecimentos.

Projetos clássicos atacam esse problema por meio de regularizações fixas (como EWC e SI) ou jogos min–max adaptativos. Neste projeto, propomos uma abordagem alternativa: modelar o equilíbrio \textbf{estabilidade–plasticidade} como um \textbf{problema de controle sequencial} resolvido via \textbf{Aprendizado por Reforço (RL)}.

A ideia é que um \textbf{agente} observe métricas do treinamento (como acurácias, perdas e deriva dos parâmetros) e \textbf{decida dinamicamente} hiperparâmetros de regularização, taxa de \textit{replay} ou congelamento de camadas, buscando maximizar o desempenho em novas tarefas sem violar limites de esquecimento.

% ======================================================
% OBJETIVOS
% ======================================================
\section{Objetivos}

\subsection*{Objetivo geral}
Investigar o uso de \textit{Reinforcement Learning} como controlador adaptativo da estabilidade–plasticidade em aprendizado contínuo, visando reduzir o esquecimento catastrófico e otimizar a retenção do conhecimento em cenários sequenciais.

\subsection*{Objetivos específicos}
\begin{enumerate}
    \item Formular o aprendizado contínuo como um problema de controle com restrição de esquecimento.
    \item Implementar um agente RL (bandit contextual e ator–crítico) para ajustar $\lambda_t$ (penalização de deriva) durante o treinamento.
    \item Comparar desempenho, estabilidade e eficiência com baselines fixos e métodos regulares (EWC, SI).
    \item Explorar variantes de controle: seleção adaptativa de taxa de \textit{replay} e congelamento seletivo de camadas.
\end{enumerate}

% ======================================================
% FORMULAÇÃO
% ======================================================
\section{Formulação como CMDP}
Modelamos o processo de aprendizado como um \textit{Constrained Markov Decision Process} (CMDP):

\paragraph{Estados.} Vetores de métricas observadas: 
\[
s_t = (\mathrm{acc}_A, \mathrm{acc}_B, \mathcal{L}_A, \mathcal{L}_B, \|\bm{\theta}-\bm{\theta}_A\|_2, \text{época}, \text{grad-norms}).
\]

\paragraph{Ações.} 
\[
a_t = (\lambda_t, \rho_t),
\]
onde $\lambda_t \ge 0$ regula a penalização de deriva e $\rho_t \in [0,1]$ é a fração de \textit{replay} opcional.

\paragraph{Recompensa e custo.}
\[
r_t = \Delta \mathrm{acc}_B, \quad c_t = \max(0,\mathrm{acc}_A^{(t-1)} - \mathrm{acc}_A^{(t)}).
\]

\paragraph{Objetivo.}
Maximizar o retorno esperado sob uma restrição de custo:
\[
\max_\pi \ \mathbb{E}_\pi\!\Big[\sum_t r_t\Big]
\quad \text{sujeito a} \quad
\mathbb{E}_\pi\!\Big[\sum_t c_t\Big] \le \tau.
\]

\paragraph{Lagrangiano.}
\[
\mathcal{L}(\pi, \lambda) = \mathbb{E}_\pi\!\Big[\sum_t (r_t - \lambda c_t)\Big], 
\quad \lambda \leftarrow [\lambda + \eta_\lambda(\bar{c} - \tau)]_+.
\]

% ======================================================
% METODOLOGIA
% ======================================================
\section{Metodologia}

\subsection*{Parte 1 — Bandit contextual (simples)}
Cada época é um episódio. O agente observa o estado $(\mathrm{acc}_A,\mathrm{acc}_B,\|\bm{\theta}-\bm{\theta}_A\|_2)$ e escolhe $\lambda_t$ de uma grade discreta. A recompensa é o ganho de $\mathrm{acc}_B$ penalizado pelo esquecimento em A. O agente é treinado via Thompson Sampling ou UCB.

\subsection*{Parte 2 — Ator–crítico com restrição}
Treinamento contínuo com ator–crítico (PPO/A2C).  
A política $\pi_\phi(a_t|s_t)$ aprende a ajustar $\lambda_t$ e $\rho_t$, enquanto o multiplicador $\lambda_\text{dual}$ é atualizado online para respeitar a restrição de esquecimento $\tau$.

\subsection*{Parte 3 — Comparações e métricas}
\begin{itemize}
    \item Baselines: $\lambda$ fixo, EWC, SI.
    \item Métricas: $\mathrm{acc}_A$, $\mathrm{acc}_B$, esquecimento total, deriva média, retorno acumulado.
    \item Gráficos: trajetória de $\lambda_t$, custo acumulado vs. limite $\tau$, curva Pareto (esquecimento × desempenho em B).
\end{itemize}

% ======================================================
% COMO RODAR
% ======================================================
\section*{Como rodar (protótipo)}
\begin{verbatim}
python rl_controller.py --epochs 10 --agent bandit --lambda_grid 1e-4 1e-3 1e-2
python rl_controller.py --agent actor_critic --theta_lr 1e-3 --policy_lr 3e-4
\end{verbatim}

% ======================================================
% EXEMPLO DE CÓDIGO
% ======================================================
\appendix
\section*{Apêndice A — Esqueleto simplificado do agente RL}

\begin{lstlisting}
# rl_controller.py
# Minimal example: bandit controlling lambda_t

for epoch in range(E):
    state = collect_metrics(model)
    lambda_t = agent.act(state)
    loss_B = CE(model(x_B), y_B)
    theta = flatten_params(model)
    F_drift = 0.5 * torch.sum((theta - theta_A)**2)
    obj = loss_B + lambda_t * F_drift
    backprop(obj)

    accA_new, accB_new = evals(...)
    reward = accB_new - accB_prev
    cost = max(0, accA_prev - accA_new)
    agent.update(state, lambda_t, reward - dual_lambda * cost)
    dual_lambda = max(0, dual_lambda + eta * (cost - tau))
\end{lstlisting}

% ======================================================
% RESULTADOS ESPERADOS
% ======================================================
\section{Resultados esperados}
Espera-se que o controlador RL:
\begin{itemize}
    \item Ajuste $\lambda_t$ dinamicamente, evitando tanto esquecimento quanto rigidez excessiva;
    \item Supere penalizações fixas em termos de estabilidade e desempenho;
    \item Produza curvas de custo e acurácia mais estáveis, aproximando-se do limite Pareto ótimo entre retenção e aprendizado.
\end{itemize}

% ======================================================
% CONCLUSÃO
% ======================================================
\section{Conclusão}
Este projeto propõe uma integração direta entre aprendizado contínuo e aprendizado por reforço. O agente atua como um controlador adaptativo, aprendendo políticas de regularização que balanceiam estabilidade e plasticidade de forma autônoma. A abordagem abre caminho para sistemas de aprendizado contínuo verdadeiramente autoajustáveis.

\end{document}
